\section{Descenso de morfismos de preesquemas}
\label{VIII.5}

\begin{proposition}
\label{VIII.5.1}
Sea $g\colon S'\to S$ un morfismo de preesquemas.
\begin{enumerate}
\item[a)] Supongamos que~$g$ es sobreyectivo y que el homomorfismo
$$
g^*\colon \mathcal{O}_{S}\to g_*(\mathcal{O}_{S'})
$$
es inyectivo. Entonces~$g$ es un epimorfismo en la categoría de los
preesquemas, e incluso en la categoría de los espacios anillados.
\item[b)] Supongamos que~$g$ es sobreyectivo y hace de~$S$ un espacio
topológico cociente de~$S'$. Sea~$S''=S'\times_{S}S'$ y sea
$h\colon S''\to S$ el morfismo estructural; consideremos el diagrama de
homomorfismos canónico:
$$
\xymatrix@C=.5cm{\mathcal{O}_{S} \ar[r] &
g_*(\mathcal{O}_{S'})\ar@/^2pt/[r]\ar@/_2pt/[r] &
h_*(\mathcal{O}_{S''})}\quoi,
$$
y supongamos que este diagrama es \emph{exacto}. Entonces~$g$ es un
epimorfismo efectivo
\index{epimorfismo efectivo|hyperpage}%
en la categoría de los preesquemas (y también en la categoría de los
espacios anillados), es decir, el diagrama
$$
\xymatrix@C=.5cm{S & S'\ar[l] & S''\ar@/^2pt/[l]\ar@/_2pt/[l]}
$$
es exacto.
\end{enumerate}
\end{proposition}

\subsubsection*{Demostración} a) Hay que demostrar que un morfismo de
espacios anillados~$f\colon S\to Z$ queda determinado cuando se conoce
$fg$. Ahora bien, como~$g$ es sobreyectivo, se conoce la aplicación entre
conjuntos~$f_{0}$ subyacente a~$f$; queda por determinar el homomorfismo
de haces de anillos~$\mathcal{O}_{Z}\to \mathcal{O}_{S}$ o, lo que es lo
mismo, el homomorfismo
\marginpar{211}
de haces de anillos
$$
u\colon f_{0}^{-1}(\mathcal{O}_{Z})\to \mathcal{O}_{S}
$$
definido por~$f$. Se conoce ya el homomorfismo
$$
(fg)_{0}^{-1}(\mathcal{O}_{Z})=
g_{0}^{-1}(f_{0}^{-1}(\mathcal{O}_{Z}))\to \mathcal{O}_{S'}
$$
definido por~$fg$ o, lo que es lo mismo, se dispone de un homomorfismo
$$
f_{0}^{-1}(\mathcal{O}_{Z})\to
g_{0*}(\mathcal{O}_{S'})=g_*(\mathcal{O}_{S'})\quoi.
$$
Se comprueba inmediatamente que este último no es sino el compuesto de
$g^*\colon \mathcal{O}_{S}\to g_*(\mathcal{O}_{S'})$ y de~$u$; como
$g^*$ es inyectivo, $u$ queda determinado cuando se conoce~$g^*u$.
[N.B.: evidentemente, no se ha utilizado que~$g\colon S'\to S$ sea un
morfismo de preesquemas; el enunciado valdría para un morfismo cualquiera
de espacios anillados. La misma observación vale para b), tanto en la
categoría de los espacios anillados como en la de los espacios localmente
anillados. Obsérvese también que, si~$g$ es un morfismo de
preesquemas no necesariamente sobreyectivo, pero tal que
$g^*\colon \mathcal{O}_{S}\to g_*(\mathcal{O}_{S'})$ sea inyectivo,
entonces, para dos morfismos~$f_{1},f_{2}$ de~$S$ en un
\emph{esquema}~$Z$ tales que~$f_{1}g=f_{2}g$, se tiene $f_{1}=f_{2}$. En
efecto, si~$I$ es el ideal sobre~$S$ que define el subpreesquema de~$S$ de
las coincidencias de~$f_{1},f_{2}$ (imagen inversa del subpreesquema
diagonal de~$Z\times Z$ por~$(f_{1},f_{2})$), se ve que~$I$ está contenido
en~$\Ker(g^*)$].

b) Hay que demostrar que, para todo espacio anillado~$Z$, el diagrama
siguiente de aplicaciones
$$
\xymatrix@C=.5cm{\Hom(S,Z)\ar[r] & \Hom(S',Z)\ar@/^2pt/[r]\ar@/_2pt/[r] &
\Hom(S'',Z)}
$$
es exacto, y que lo mismo sucede cuando~$Z$ es un espacio localmente
anillado y uno se limita a los homomorfismos de espacios localmente
anillados. Como ya se sabe por a) que la primera aplicación es
inyectiva, queda por ver que, si~$f'\colon S'\to Z$ es un homomorfismo de
espacios anillados tal que~$f'p_{1}=f'p_{2}$,
\marginpar{212}
entonces~$f'$ es de la forma~$fg$, donde~$f\colon S\to Z$ es un
homomorfismo de espacios anillados. Como~$g$ es sobreyectivo, resulta
entonces evidente que, si~$f'$ es un morfismo de espacios localmente
anillados, lo mismo sucederá con~$f$.

De la hipótesis sobre~$f'$ resulta que la aplicación entre conjuntos
subyacente~$f_{0}'$ es constante sobre las fibras de la
aplicación~$g_{0}$; como esta última es sobreyectiva, $f_{0}'$ se
factoriza de manera única como~$f_{0}'=f_{0}g_{0}$, donde
$f_{0}\colon S\to Z$ es una aplicación, necesariamente continua porque
$g_{0}$ identifica~$S$ con un espacio topológico cociente de~$S'$.
Consideremos ahora el homomorfismo
$$
f_{0}^{-1}(\mathcal{O}_{Z})\to g_*(\mathcal{O}_{S'})
$$
deducido del homomorfismo~$(f_{0}g_{0})^{-1}(\mathcal{O}_{Z})\to
\mathcal{O}_{S'}$ correspondiente a~$f'$. La hipótesis
$f'p_{1}=f'p_{2}$ se interpreta entonces diciendo que los compuestos del
homomorfismo precedente con los dos homomorfismos
$\xymatrix@C=.5cm{g_*(\mathcal{O}_{S'})\ar@/^2pt/[r]\ar@/_2pt/[r] &
h_*(\mathcal{O}_{S''})}$ son iguales; por la hipótesis b), aquel se
factoriza, por tanto, a través de un morfismo
$$
f_{0}^{-1}(\mathcal{O}_{Z})\to \mathcal{O}_{S}\quoi.
$$
Este último define un morfismo de espacios anillados~$f\colon S\to Z$,
que es el morfismo buscado.

\begin{theoreme}
\label{VIII.5.2}
Sea $\cal{F}$ la categoría fibrada de las flechas en la categoría~$\Sch$
de los preesquemas \textup{(\Ref{VI.11}.a)}. Entonces todo
morfismo~$g\colon S'\to S$ fielmente plano y cuasicompacto es un morfismo
de descenso para~$\cal{F}$ (o también, como se dice, un morfismo de
descenso en~$\Sch$).
\index{descenso (morfismo de)|hyperpage}%
\end{theoreme}

Esto significa lo siguiente: sea~$S''=S'\times_{S}S'$ y, para dos
preesquemas~$X,Y$ sobre~$S$, consideremos sus imágenes inversas~$X',Y'$ sobre
$S'$ y sus imágenes inversas~$X'',Y''$ sobre~$S''$; se obtiene un diagrama
de aplicaciones
$$
\xymatrix@C=.5cm{\Hom_{S}(X,Y)\ar[r] &
\Hom_{S'}(X',Y')\ar@/^2pt/[r]\ar@/_2pt/[r] & \Hom_{S''}(X'',Y'')}\quoi;
$$
sentadas estas notaciones, \Ref{VIII.5.2} significa que este diagrama es
exacto. Obsérvese que
\marginpar{213}
no es cierto, en general, que~$g$ sea un morfismo de descenso efectivo,
es decir, que para todo preesquema~$X'$ sobre~$S'$, todo dato de descenso
sobre~$X'$ con respecto a~$g\colon S'\to S$ sea efectivo. La cuestión de
la efectividad, a menudo delicada, se examinará en el
\No~\Ref{VIII.7}.

Se ha visto en~\cite{VIII.D} (teniendo en cuenta que en~$\Sch$ existen los
productos fibrados) que el enunciado~\Ref{VIII.5.2} equivale al siguiente:

\begin{corollaire}
\label{VIII.5.3}
Todo morfismo fielmente plano y cuasicompacto de preesquemas es un
epimorfismo efectivo universal.
\end{corollaire}

Como un morfismo fielmente plano y cuasicompacto sigue siéndolo tras toda
extensión de la base, se reduce uno a demostrar que es un epimorfismo
efectivo. Se aplica entonces el criterio~\Ref{VIII.5.1} b), que da el
resultado buscado teniendo en cuenta~\Ref{VIII.4.3} y~\Ref{VIII.1.7}.

\begin{corollaire}
\label{VIII.5.4}
Sea $g\colon S'\to S$ un morfismo fielmente plano y cuasicompacto,
$f\colon X\to Y$ un~$S$-morfismo y $f'\colon X'\to Y'$ el
$S'$-morfismo que se deduce de él por el cambio de base~$S'\to S$. Para
que~$f$ sea un isomorfismo, es necesario y suficiente que~$f'$ lo sea.
\end{corollaire}

En efecto, si $f'$ es un isomorfismo, es también un isomorfismo para las
estructuras de descenso naturales sobre~$X',Y'$; como el
funtor~$X\mto X'$ de~$\Sch_{/S}$ en la categoría de los preesquemas
sobre~$S'$ provistos de un dato de descenso con respecto a~$g$ es plenamente
fiel por~\Ref{VIII.5.2}, se obtiene la conclusión buscada.

\begin{corollaire}
\label{VIII.5.5}
Bajo las condiciones de~\Ref{VIII.5.4}, para que~$f$ sea una inmersión
cerrada (\resp una inmersión abierta, \resp una inmersión cuasicompacta),
\ifthenelse{\boolean{orig}}
{es necesario y suficiente}
{es necesario y suficiente}
que~$f'$ lo sea.
\end{corollaire}

Puede suponerse, como de costumbre, $Y=S$, y solo hay que demostrar la
suficiencia. Obsérvese que el hecho de que~$X'/Y'$ esté provisto de un dato
de descenso con respecto a~$g\colon Y'\to Y$ y que el morfismo
estructural~$f'\colon X'\to Y'$ sea una inmersión y, por tanto, un
monomorfismo, implica que los dos subobjetos de~$Y''$ obtenidos como
imágenes
\marginpar{214}
inversas de~$X'/Y'$ por una y otra proyección de~$S''$ en~$S'$ son los
mismos. Si~$f'$ es una inmersión cerrada, de~\Ref{VIII.1.9} resulta que
existe un subpreesquema cerrado~$X_{1}$ de~$Y$ cuya imagen inversa por
$g\colon Y'\to Y$ es~$X'$. Por la unicidad de la solución de un problema
de descenso con respecto a un morfismo de descenso para~$\cal{F}$,
\ifthenelse{\boolean{orig}}
{}
{se}
sigue que~$X_{1}$ es $Y$-isomorfo a~$X$; por tanto,
$f\colon X\to Y$ es una inmersión cerrada. Se procede del mismo modo para
una inmersión abierta, utilizando~\Ref{VIII.4.4}. Finalmente, si~$f'$ es
una inmersión cuasicompacta, $f$ es cuasicompacto en virtud
de~\Ref{VIII.3.3}; puede aplicarse entonces al subconjunto~$f(X)$ de~$Y$
el criterio~\Ref{VIII.4.5}, que demuestra que~$f(X)$ es localmente cerrado,
pues su imagen inversa~$f'(X')$ en~$Y'$ lo es. Reemplazando entonces~$Y$
por un subconjunto abierto en el que~$f(X)$ sea cerrado, se reduce uno al
caso en que~$f'$ es una inmersión cerrada; por lo anterior, también lo
es~$f$.

\begin{corollaire}
\label{VIII.5.6}
Bajo las condiciones de~\Ref{VIII.5.4}, para que~$f$ sea afín, es necesario
y suficiente que~$f'$ lo sea.
\end{corollaire}

Se procede como en~\Ref{VIII.5.5}, utilizando~\Ref{VIII.2.1} (también
puede utilizarse el criterio cohomológico de Serre [EGA II 5.2], que
demuestra~\Ref{VIII.5.6} sin utilizar ninguna técnica de descenso).

\begin{corollaire}
\label{VIII.5.7}
Bajo las condiciones de~\Ref{VIII.5.4}, para que~$f$ sea entero (\resp
finito, \resp finito y localmente libre), es necesario y suficiente
que~$f'$ lo sea.
\end{corollaire}

Solo hay que demostrar la suficiencia y, como de costumbre, puede
suponerse $Y=S$, con $Y$ afín e~$Y'$ afín. Como la hipótesis implica que
$f'$ es afín, lo mismo sucede con~$f$ por~\Ref{VIII.5.6}; por tanto, $X$
y, en consecuencia, $X'$ son afines. Sean~$A,A',B,B'=B\otimes_{A}A'$ los
anillos de~$Y,Y',X,X'$. Se tiene $B=\varinjlim_i B_{i}$, donde~$B_{i}$
recorre las sub-$A$-álgebras de~$B$ que son de tipo finito sobre~$A$; de
aquí $B'=\varinjlim_i B_{i}'$, donde las~$B_{i}'$ son subálgebras de tipo
finito de la $A'$-álgebra~$B'$. Si $B'$ es entero
\ifthenelse{\boolean{orig}}
{sobre~$A$,}
{sobre~$A'$,}
las~$B_{i}'$ son módulos de tipo finito sobre~$A'$; como $A'$ es fielmente
plano sobre~$A$, las~$B_{i}$ son módulos de tipo finito sobre~$A$, es decir,
$B$ es entero sobre~$A$. Se ve del mismo modo que, si~$B'$ es finito
sobre~$A'$, $B$ lo es sobre~$A$. La misma conclusión vale para «localmente
libre de tipo finito»; \cf \Ref{VIII.1.11}.

\begin{corollaire}
\label{VIII.5.8}
Bajo
\marginpar{215}
las condiciones de~\Ref{VIII.5.4}, supongamos~$f$ cuasicompacto y sean
$\mathcal{L}$ un haz invertible sobre~$X$ y~$\mathcal{L}'$ su imagen
inversa sobre~$X'$. Para que~$\mathcal{L}$ sea amplio (\resp muy amplio)
con respecto a~$f$, es necesario y suficiente
\ifthenelse{\boolean{orig}}
{que~$f'$}
{que~$\cal{L}'$}
sea amplio (\resp muy amplio) con respecto a~$f'$.
\end{corollaire}

Solo hay que demostrar la suficiencia. La hipótesis sobre
\ifthenelse{\boolean{orig}}
{$\mathcal{L}$}
{$\mathcal{L}'$}
implica, en todo caso, que~$f'$ es separado; por~\Ref{VIII.4.8}, $f$ es
separado. Como~$f$ es cuasicompacto y~$g\colon Y'\to Y$ es plano, el
cálculo de las imágenes directas mediante recubrimientos afines demuestra
que se tienen isomorfismos
$$
g^*(f_*(\mathcal{L}^{\otimes n}))\isomto
f_*'(\mathcal{L}'^{\otimes n})
$$
para todo entero~$n$ y, por tanto, un isomorfismo
$$
g^*(\mathcal{S})\isomto \mathcal{S}'\quoi,
$$
donde~$\mathcal{S}$ (\resp $\mathcal{S}'$) denota el álgebra graduada
cuasicoherente sobre~$Y$ (\resp sobre~$Y'$) suma directa de
los~$f_*(\mathcal{L}^{\otimes n})$ (\resp de
los~$f_*'(\mathcal{L}'^{\otimes n})$) para~$n\geq 0$. Obsérvese que, para
todo~$n\geq 0$, el conúcleo del homomorfismo
canónico~$f'^*(\mathcal{S}_{n}')\to \mathcal{L}'^{\otimes n}$ es la
imagen inversa por~$X'\to X$ del conúcleo de
$f^*(\mathcal{S}_{n})\to \mathcal{L}^{\otimes n}$; por tanto, su
soporte~$Z_{n}'$ es la imagen inversa del soporte~$Z_{n}$. Si
\ifthenelse{\boolean{orig}}
{$f$}
{$\cal{L}'$}
es amplio, la intersección de los~$Z_{n}'$ es vacía; como~$X'\to X$ es
sobreyectivo, la intersección de los~$Z_{n}$ es vacía, es decir, se tiene
un morfismo canónico
$$
j\colon X\to \Proj(\mathcal{S})
$$
(EGA~II 3). Además, el morfismo análogo
$$
j'\colon X'\to \Proj(\mathcal{S}')
$$
no es sino el que se deduce del precedente mediante el cambio de
base~$Y'\to Y$ (\loccit). Sentado esto, decir que~$\mathcal{L}'$ es amplio
con respecto a~$f'$ significa que~$j'$ es una inmersión, necesariamente
cuasicompacta, además, porque~$f'$ es cuasicompacto. Por
\Ref{VIII.5.5}, $j$ es, por tanto, una inmersión, es decir,
$\mathcal{L}$ es amplio con respecto a~$f$. --- Se procede de manera
enteramente análoga en el caso «muy
\marginpar{216}
amplio», limitándose arriba a~$n=1$ y reemplazando la consideración de
$\Proj(\mathcal{S})$ por la del fibrado proyectivo
$\mathcal{P}(\mathcal{S}_{1})$ asociado a~$\mathcal{S}_{1}$.

Recordemos (EGA~II 5.1.1) que un morfismo cuasicompacto~$f$ se llama
\emph{cuasiafín} si, para todo abierto afín~$U$ de~$Y$, $f^{-1}(U)$ es un
preesquema isomorfo a un subesquema abierto de un esquema afín. Se
demuestra (\loccit) que equivale a decir que~$\mathcal{O}_{X}$ es amplio
(o también muy amplio) con respecto a~$f$. Por tanto, \Ref{VIII.5.8}
implica:

\begin{corollaire}
\label{VIII.5.9}
Bajo las condiciones de~\Ref{VIII.5.4} y suponiendo~$f$ cuasicompacto,
para que~$f$ sea cuasiafín, es necesario y suficiente que~$f'$ lo sea.
\end{corollaire}

\begin{remarques}
\label{VIII.5.10}
El ejemplo de variedad no proyectiva de Hironaka muestra que puede tenerse
un morfismo propio~$f\colon X\to Y$ de variedades algebraicas no singulares
(con~$Y$ proyectiva), tal que~$Y$ sea reunión de dos abiertos~$Y_{i}$ para
los cuales~$X_{i}=X\times_{Y}Y_{i}$ sea proyectivo sobre~$Y_{i}$, pero sin
que~$f$ sea proyectivo. Poniendo entonces $Y'=Y_{1}\amalg Y_{2}$, $Y'$ es
fielmente plano y cuasicompacto (e incluso cuasifinito) sobre~$Y$,
$f'\colon X'\to Y'$ es proyectivo, pero~$f$ no lo es. Hay que tener
cuidado, por tanto: para aplicar~\Ref{VIII.5.8} y deducir del hecho de
que~$f'$ sea proyectivo la misma conclusión sobre~$f$, hay que disponer ya
sobre~$X'$ de un haz invertible~$\mathcal{L}'$ amplio para~$f'$,
\emph{provisto de un dato de descenso con respecto a}~$X'\to X$ (lo que
permite considerar~$\mathcal{L}'$ como la imagen inversa de un haz
invertible~$\mathcal{L}$ sobre~$X$, que será entonces amplio para~$f$
gracias a~\Ref{VIII.5.8}). No obstante, cuando~$g\colon S'\to S$ es finito
y localmente libre, véase~\Ref{VIII.7.7}.
\end{remarques}
